\documentclass[10pt,letterpaper]{article}
\usepackage[margin=0.7in,headheight=14pt]{geometry}
\usepackage[T1]{fontenc}
\usepackage{lmodern,amsmath,booktabs,xcolor,fancyhdr,microtype}
\usepackage[colorlinks=true,linkcolor=teal,urlcolor=teal]{hyperref}
\definecolor{navy}{HTML}{192B4D}
\definecolor{teal}{HTML}{007D80}
\definecolor{pale}{HTML}{F3F6F8}
\pagestyle{fancy}\fancyhf{}
\fancyhead[L]{\small MF 731 | Homework 3}
\fancyhead[R]{\small Ten-Day VaR}
\fancyfoot[C]{\small\thepage}
\setlength{\parindent}{0pt}
\setlength{\parskip}{6pt}
\hypersetup{pdftitle={MF 731 Homework 3 - Ten-Day VaR}}
\begin{document}
{\LARGE\bfseries\textcolor{navy}{Ten-Day VaR of Delta-Hedged Puts}}

\textbf{Setup.} I use the assigned rounded parameters:
$S_0=152.51$, $K=170$, $T=0.292$, $r=0.00119$,
$\sigma=0.4907$, $\mu=0.1691$ and $\Delta=1/252$.
There are 100 short puts, with no additional contract multiplier.
$P$ is the Black--Scholes \emph{put} price; ``call'' in the handout is a typo.
Loss is $L_j=100(V_0-V_j)$, where $V$ is value per put.
I keep the initial stock hedge fixed over both horizons and let the initial
cash balance earn interest, using the handout's terminal-value formula.

\textbf{1. One-day simulation and time scaling.}
For each of 40,000 paths, simulate ten independent standard normal shocks:
\[
S_j=S_{j-1}\exp\!\left[
  \left(\mu-\tfrac12\sigma^2\right)\Delta
  +\sigma\sqrt{\Delta}\,Z_j\right].
\]
I use seed 731 and the first day of these paths for the one-day calculation.
Initially $h_0=\partial_sP(0,S_0)=-0.608555$ and
$Y_0=-h_0S_0$, as specified, so $V_0=-P(0,S_0)$.
The first-day value is
\[
V_1=h_0S_1+Y_0e^{r\Delta}-P(\Delta,S_1).
\]
The empirical 95\% VaR is the 38,000th sorted loss, with no interpolation.
This gives a one-day VaR of \textbf{\$30.14} and
\[
\widehat{\mathrm{VaR}}_{10}^{\mathrm{scale}}
=\sqrt{10}\,\widehat{\mathrm{VaR}}_1
=\boxed{\$95.32}.
\]

\textbf{2. Ten-day time aggregation.}
Aggregate the ten simulated daily log returns within each path:
\[
X_{10}=\sum_{j=1}^{10}\log(S_j/S_{j-1}),
\qquad S_{10}=S_0e^{X_{10}}.
\]
Then reprice the put with remaining maturity $T-10\Delta$ and calculate
\[
V_{10}=h_0S_{10}+Y_0e^{10r\Delta}-P(10\Delta,S_{10}).
\]
The 95th percentile of $100(V_0-V_{10})$ is
\[
\widehat{\mathrm{VaR}}_{10}^{\mathrm{aggregate}}=\boxed{\$308.28}.
\]

\textbf{3. Comparison.}
Time aggregation gives a VaR \$212.96 higher than time scaling,
or about 3.23 times as large. The square-root rule is exact for sums of
independent zero-mean normal losses. Here the initial delta hedge
cancels the linear price term, leaving mainly the short put's gamma exposure.
The leading random loss term is quadratic in the stock move. Ten-day log
returns have ten times the one-day variance, so their squared shocks scale
roughly with ten rather than $\sqrt{10}$. The fixed hedge also becomes less
effective after a large move. These effects explain why time scaling
substantially understates the fully repriced ten-day VaR.

\vfill
\begin{center}
\colorbox{pale}{\parbox{0.94\linewidth}{\centering
\textbf{Python code and running instructions}\\[4pt]
\url{https://github.com/jesuswalking/MF-731-HW1/tree/main/HW3}
}}
\end{center}
\end{document}
